% TOP_PROBLEM: {"id": 3000042, "title": "Highly element-connected orientation", "category_id": 2, "category": {"id": 2, "name": "combinatorics", "display_name": "Combinatorics", "description": "Counting problems, graph theory, discrete structures.", "slug": "combinatorics", "order_index": 2, "created_at": "2026-07-31T15:26:25.671Z"}, "status": "open", "problem_number": "AMR-029-0042", "set_id": 13}
% FIRST_POSTED: 2026-10-02
% ENTRY_KIND: statement_only
% UnsolvedMath ID 3000042; source catalogue code AMR-029-0042
% !TeX program = lualatex
% Definition and sources only; this file is not an LLM solution attempt.
\documentclass[11pt,a4paper]{article}
\usepackage[margin=25mm]{geometry}
\usepackage{fontspec}
\setmainfont{lmroman10-regular.otf}[BoldFont=lmroman10-bold.otf,
 ItalicFont=lmroman10-italic.otf,BoldItalicFont=lmroman10-bolditalic.otf]
\usepackage{amsmath,amssymb,mathtools}
\usepackage{unicode-math}
\setmathfont{latinmodern-math.otf}
\AtBeginDocument{\let\setminus\smallsetminus}
\usepackage{xurl,hyperref}
\hypersetup{colorlinks=true,urlcolor=blue}
\setcounter{secnumdepth}{0}
\setlength{\parindent}{0pt}
\setlength{\parskip}{5pt}
\setlength{\emergencystretch}{3em}
\begin{document}
\section*{Highly element-connected orientation}\label{3000042}
{\small UnsolvedMath ID 3000042 · AMR-029-0042 · Combinatorics · AMR Open Problem Lists}
\subsection{Definitions and mathematical statement}
Let $G=(V,E)$ be a finite loopless undirected graph and let
$T\subseteq V$ be its terminal set. Vertices in $V\setminus T$ are
Steiner vertices. The elements that may fail are the edges and the
Steiner vertices; terminals are not elements.

Paths are element-disjoint when they have no common edge and no
common Steiner vertex. They may meet at terminals, including terminals
other than their endpoints. The graph is $q$-element-connected on $T$
if every pair of distinct terminals is joined by $q$ such paths.
Equivalently, deleting fewer than $q$ elements cannot disconnect any
two terminals.

For every integer $k\ge1$, is the following assertion true?
\[
 G\text{ is }2k\text{-element-connected on }T
 \ \Longrightarrow\
 G\text{ has a }k\text{-element-connected orientation on }T.
\]
An orientation chooses a direction for each edge. Its required property
means that for every ordered pair of distinct terminals $s,t$, there
are $k$ directed element-disjoint paths from $s$ to $t$.

The same orientation must work for all ordered terminal pairs.
Paths used for different pairs need not be disjoint from each other;
disjointness applies within each family of $k$ paths. Choosing a root
and satisfying only the paths from that root is a weaker target.
Likewise, ordinary edge-connectivity ignores failure of Steiner vertices
and cannot simply replace element-connectivity in the hypothesis.
Parallel edges can be retained as distinct elements.
\subsection{Short English statement}
Does twice the required terminal element-connectivity guarantee an orientation preserving the required directed connectivity between every pair of terminals?
\subsection{Sources}
\begin{itemize}
\item[S1] ULAM.AI and the UnsolvedMath Contributors, supplied problems.json, record 3000042 (AMR-029-0042), AMR Open Problem Lists. Catalogue identity and source transcription; CC BY 4.0.
\url{https://huggingface.co/datasets/ulamai/UnsolvedMath}.
\item[S2] Egres Open - List of open problems, Open-problem page: Highly element-connected orientation. Original source indexed by AMR-029-0042.
\url{https://oldlemon.cs.elte.hu/egres/open/List_of_open_problems}.
\item[S3] EGRES Open, Highly element-connected orientation. Individual source statement and remarks.
\url{https://lemon.cs.elte.hu/egres/open/Highly_element-connected_orientation}.
\end{itemize}
\end{document}
